\scanpage{218}
\mspage{134}
% archive sheet number: 217
% element: text | top, lines 1-3 | Enfin revenant a la dem
\deleted{La dém.} Enfin (revenant à la dém. de 4.5.13, 3\textsuperscript{o}), l'implication $\unsure{3|B})\Rightarrow1)$ provient du fait que $M_\pi$ est stable par $\varinjlim$ de cardinal $\leq\pi$.

% element: text | upper third, lines 4-9 | Reste a prouver
Reste à prouver que ces conditions, \unsure{ici|i.e.} la stabilité de $\mathrm{Filt}'_\pi$ par composition de autres flèches, est satisfaite dans les cas A) et C), le cas A) étant d'ailleurs déjà connu. Le cas B) résulte \deleted{\unsure{facile|facilement}} immédiatement de la transitivité des familles épimorphiques strictes universelles, \deleted{\unsure{Je crois que le cas}} et c'est d'ailleurs \unsure{écrit en toutes} explicitement \deleted{\unsure{dans}} (la stabilité \inserted{de $\mathrm{Filt}'{}^\pi$} par $\varinjlim$ de cardinal $\leq\pi$) dans SGA 4 I 9.13. (\unsure{Non pas} dans l'énoncé, mais c'est rappelé dans la démonstration).

% element: text | middle lower, recollection of SGA | En redigeant
En rédigeant, \inserted{SGA 4 I 9} je savais donc que $M_\pi\subset\mathrm{Filt}'{}^\pi$ pour $\forall\pi$, mais l'inclusion, plus évidente pour \unsure{tout|tt} \inserted{(dans le cas $\unsure{\pi=\pi_0|\pi\geq\pi_0}$)}, $\mathrm{Filt}'{}^\pi\subset M_\pi$ (quand on suppose les objets de $C$ $\unsure{\pi_0|\pi_p}$-accessibles) m'avait apparemment échappé !

% element: text | lower third, tentative comparison | Reste a traiter le cas C
Reste à traiter le cas C. J'ai l'impression qu'il est \inserted{un peu} plus général que le cas B), i.e. que \inserted{la condition sur $M$} B implique C, mais non inversement. Ainsi la condition C est vérifiée dès que $M$ \unsure{correspond} : une espèce de \unsure{structures|structure} alg. \unsure{discrète|abstraite} \deleted{\unsure{par|pour}} \inserted{un}
